\section{Validasi \textit{Barren Plateau}}
Keberhasilan implementasi algoritma hibrida kuantum-klasik pada era \textit{Noisy Intermediate-Scale Quantum} (NISQ) sangat bergantung pada efektivitas proses optimasi parameter dalam menavigasi lanskap energi yang kompleks \cite{preskill_quantum_2018}. Namun, kemunculan fenomena patologis seperti \textit{Barren Plateau} sering kali menghambat aliran informasi gradien sehingga menyebabkan stagnasi pada proses pencarian solusi optimal di ruang Hilbert \cite{mcclean_barren_2018}. Validasi terhadap kondisi sistem selama fase optimasi menjadi prasyarat krusial guna menjamin bahwa lintasan stokastik parameter sirkuit tetap berada pada zona yang dapat dilatih secara efektif. Bab ini akan menguraikan mekanisme diagnostik melalui pemanfaatan metrik entropi Von Neumann dan analisis statistik varians gradien guna mendeteksi serta memitigasi hambatan komputasional tersebut \cite{navarrete_quantum_2022}. Pemahaman mendalam mengenai limitasi fisik ini akan memberikan landasan teoretis bagi penggunaan strategi inisialisasi cerdas berbasis teori permainan guna meningkatkan stabilitas konvergensi portofolio.

\subsection{\textit{Barren Plateau}}
Fenomena \textit{Barren Plateau} atau dataran tandus didefinisikan sebagai kondisi patologis di mana lanskap fungsi biaya pada algoritma kuantum variasional menjadi sangat datar secara eksponensial seiring bertambahnya jumlah \textit{qubit} di dalam sistem. Kondisi tersebut dicirikan oleh hilangnya sinyal gradien secara masif, sehingga algoritma optimasi berbasis gradien kehilangan kemampuan fungsionalnya dalam menentukan arah pergerakan menuju konfigurasi energi minimum yang diinginkan. Secara fisik, kemunculan dataran tandus ini sering kali dipicu oleh inisialisasi parameter yang terlalu acak di dalam sirkuit kuantum yang memiliki derajat keterbelitan (\textit{entanglement}) yang sangat tinggi \cite{mcclean_barren_2018}. Masalah skalabilitas ini menjadi hambatan operasional utama dalam implementasi perangkat keras era NISQ karena membatasi kedalaman sirkuit yang dapat dilatih secara efektif tanpa mengalami stagnasi proses pembelajaran. Deteksi dini terhadap kemunculan dataran tandus menjadi sangat krusial guna menjamin efisiensi komputasi serta keberhasilan konvergensi pada sistem kuantum yang memiliki dimensi ruang Hilbert berskala besar.

Dampak dari fenomena \textit{Barren Plateau} sangat merugikan bagi efisiensi komputasi karena menuntut jumlah pengukuran (\textit{shots}) yang sangat besar hanya untuk mengekstraksi sinyal gradien dari dominasi derau perangkat keras secara akurat. Ketika varians dari gradien tersebut menyusut di bawah ambang batas presisi alat ukur, sirkuit kuantum secara efektif akan kehilangan kapasitas fungsionalnya untuk melakukan pembaruan parameter yang bermakna bagi proses konvergensi. Fenomena stagnasi ini memaksa komunitas peneliti untuk mengembangkan berbagai teknik inisialisasi cerdas guna menjaga agar sinyal gradien tetap berada pada level yang dapat terdeteksi selama fase optimasi berlangsung \cite{mcclean_barren_2018}. Dalam penelitian ini, integrasi teori permainan melalui \textit{Exact Potential Game} (EPG) diusulkan sebagai solusi preventif guna memberikan inisialisasi parameter yang lebih terarah dan keluar dari zona tandus tersebut. Dengan memanfaatkan titik ekuilibrium Nash sebagai konfigurasi awal, algoritma VQE dapat melewati fase pencarian acak yang rentan terhadap degradasi gradien di dalam ruang Hilbert yang sangat luas.

\subsection{Entropi Von Neumann}
Entropi Von Neumann dimanfaatkan sebagai metrik kuantitatif yang sangat esensial untuk mengukur derajat keterbelitan (\textit{entanglement}) di dalam sistem kuantum multi-qubit yang sedang menjalani proses optimasi hibrida secara iteratif dan berkelanjutan. Konsep ini memperluas paradigma entropi Shannon guna mengevaluasi derajat kemurnian suatu keadaan kuantum melalui operator densitas $\rho$ yang secara formal dinyatakan dalam bentuk persamaan matematis berindeks sebagai berikut:
\begin{equation}
S(\rho) = -\text{Tr}(\rho \log \rho)
\end{equation}
Dalam arsitektur \textit{Parameterized Quantum Circuits} (PQC), peningkatan derajat keterbelitan yang tidak terkontrol terbukti berkorelasi kuat dengan hilangnya lanskap gradien secara sistemik di seluruh ruang parameter sirkuit yang sedang dioptimasi \cite{mcclean_barren_2018}. Secara teoretis, sirkuit yang telah mencapai karakteristik \textit{2-design} cenderung memiliki nilai entropi yang mendekati batas maksimum atau yang secara akademis didefinisikan sebagai \textit{Page entropy} \cite{navarrete_quantum_2022}. Nilai entropi yang sangat tinggi tersebut menandakan bahwa informasi kuantum telah tersebar secara merata di seluruh ruang Hilbert sehingga menyebabkan fungsi biaya menjadi sulit untuk dieksplorasi oleh algoritma optimasi klasik.

Nilai numerik dari entropi Von Neumann memberikan indikasi fisik yang sangat langsung mengenai status jalinan sistem, di mana nilai $S=1$ secara definitif mengonfirmasi terjadinya keterbelitan maksimal yang merepresentasikan Keadaan Bell. Keberadaan nilai entropi pada rentang $0 < S < 1$ menunjukkan kondisi keterbelitan parsial yang menandakan adanya korelasi kuantum antar-\textit{qubit} namun belum mencapai titik jenuh informasi sistemik yang melumpuhkan sinyal gradien. Sebaliknya, nilai $S=0$ memberikan jaminan fisis bahwa sistem berada dalam keadaan terpisahkan (\textit{separable state}) tanpa adanya korelasi kuantum antar-sub-sistem yang dapat dideteksi melalui mekanisme pengukuran operator densitas \cite{nielsen_quantum_2010}. Pemanfaatan entropi sebagai instrumen diagnostik strategis memungkinkan peneliti untuk memantau apakah sirkuit sedang mendekati ambang batas \textit{2-design} yang secara teoretis memicu munculnya fenomena \textit{Barren Plateau} yang merugikan. Dengan demikian, metrik entropi ini berfungsi sebagai parameter kontrol krusial untuk menjaga agar sistem tetap berada dalam zona yang dapat dioptimasi secara efektif guna mencapai target konvergensi solusi.

Dinamika fluktuasi nilai entropi Von Neumann selama proses iterasi hibrida mencerminkan evolusi struktural pada vektor keadaan sistem yang dipicu oleh pembaruan parameter oleh unit pemrosesan klasik secara kontinu dan sistematis. Penurunan nilai entropi dari kondisi maksimal menuju titik deterministik menandakan transisi sistem dari status keterbelitan acak menuju konvergensi solusi portofolio yang bersifat terpisahkan pada tahap akhir proses optimasi sistem. Sebaliknya, peningkatan nilai entropi selama tahapan eksplorasi menunjukkan bahwa sirkuit sedang menavigasi ruang Hilbert guna meminimalkan fungsi biaya yang mengakibatkan distribusi probabilitas pengukuran menjadi tersebar secara merata \cite{navarrete_quantum_2022}. Fenomena tersebut secara empiris membuktikan bahwa derajat keterbelitan di dalam sirkuit PQC bersifat sangat dinamis dan mampu bertransformasi sebagai respons terhadap pencarian titik minimum pada lanskap energi portofolio finansial. Observasi sistematis terhadap evolusi entropi ini memberikan wawasan mendalam mengenai mekanisme internal sistem kuantum dalam memproses informasi korelasi aset selama berlangsungnya siklus optimasi hibrida menuju titik keseimbangan optimal.

\subsection{Varians Gradien}
Fenomena \textit{Barren Plateau} secara formal dapat didefinisikan melalui analisis terhadap perilaku statistik dari sinyal gradien fungsi biaya terhadap seluruh parameter sirkuit yang dioptimasi secara iteratif pada unit pemrosesan klasik. Pada kondisi stagnasi ini, nilai rata-rata dari vektor gradien cenderung mendekati titik nol mutlak, namun parameter statistik yang paling krusial adalah nilai varians gradien yang menghilang secara eksponensial \cite{mcclean_barren_2018}. Secara matematis, penyusutan varians gradien $\text{Var}[\partial_{\theta} E]$ akan mengikuti pola peluruhan asimtotik $\mathcal{O}(2^{-n})$ di mana variabel $n$ merepresentasikan jumlah unit \textit{qubit} yang dioperasikan di dalam arsitektur komputer kuantum sebagai berikut:
\begin{equation}
\text{Var}[\partial_{\theta} E] \approx \mathcal{O}(2^{-n})
\end{equation}
Penyusutan varians yang bersifat sangat masif ini mengakibatkan sinyal gradien yang bermakna terkubur sepenuhnya di bawah level derau perangkat keras (\textit{hardware noise}) serta batasan presisi numerik yang tersedia pada sistem. Oleh karena itu, mitigasi terhadap pelemahan varians gradien menjadi prasyarat mutlak bagi keberhasilan implementasi algoritma kuantum pada masalah finansial yang melibatkan jumlah aset investasi dalam skala yang sangat signifikan.
